\documentclass[11pt]{article}

\usepackage[a4paper,margin=1in]{geometry}
\usepackage{amsmath,amssymb,amsthm}
\usepackage{enumitem}
\usepackage{xcolor}
\usepackage{fancyhdr}
\usepackage{needspace}
\usepackage{physics}

\definecolor{courseblue}{RGB}{0,0,255}

\setlength{\parindent}{0pt}
\setlength{\parskip}{0.35em}
\setlength{\headheight}{14pt}

\pagestyle{fancy}
\fancyhf{}
\fancyfoot[C]{\thepage}
\renewcommand{\headrulewidth}{0pt}

\renewcommand{\ip}[2]{\left\langle #1, #2 \right\rangle}

\usepackage{bm}
\renewcommand{\vec}[1]{\bm{#1}} % Converts all \vec{} commands to bold

\DeclareMathOperator{\dist}{dist}
\DeclareMathOperator{\open}{open}
\DeclareMathOperator{\closed}{closed}
\newtheorem*{lemma}{Lemma}
\newcounter{problem}
\newenvironment{problem}{%
    \newpage
    \refstepcounter{problem}%
    \Needspace{7\baselineskip}%
    \par\vspace{0.7em}%
    \noindent\textcolor{courseblue}{\rule{\textwidth}{0.4pt}}%
    \par\vspace{0.5em}%
    \begin{list}{}{%
        \setlength{\leftmargin}{2.3em}%
        \setlength{\labelwidth}{1.8em}%
        \setlength{\labelsep}{0.5em}%
        \setlength{\topsep}{0pt}%
        \setlength{\partopsep}{0pt}%
    }%
    \item[{[\theproblem]}]
}{%
    \end{list}
    \par\vspace{0.5em}
}

\newenvironment{parts}{%
    \begin{enumerate}[label={[\alph*]},leftmargin=2.2em,itemsep=0.65em,topsep=0.65em]
}{%
    \end{enumerate}
}

% Type each solution between \begin{answer} and \end{answer}.
\newenvironment{answer}{%
    \Needspace{4\baselineskip}%
    \par\vspace{0.35em}
    \noindent\textbf{Answer:}\par
}{%
    \par\vspace{0.25em}
}

\begin{document}

\begin{center}
    {\color{courseblue}\Large\textbf{AMATH/PMATH 331: Applied Real Analysis}}
    \hfill
    {\color{courseblue}\Large\textbf{Fall 2026}}

    \vspace{1.2em}
    {\color{courseblue}\Large
    Assignment 3a Submission; due Tuesday, October 20, 2026}

    \vspace{1.2em}
    
    \textbf{Name:} Vlad Paun

    \textbf{Student number:} 21291453
\end{center}

\begin{problem}
Let $f:A\subseteq\mathbb{R}^n\to\mathbb{R}^m$.
\begin{parts}
    \item Let $B\subseteq\mathbb{R}^m$. Show that
    $f^{-1}(B^c)=A\cap(f^{-1}(B))^c$.
    \begin{answer}
        \begin{proof}
            Take arbitrary $x\in f^{-1}(B^c)$.
            \begin{align*}
                & \iff x\in A \wedge f(x)\in B^c \\
                & \iff x\in A \wedge f(x)\not\in B\\
                & \iff x\in A \wedge x\not\in f^{-1}(B)\\
                & \iff x\in A \wedge x\in (f^{-1}(B))^c\\
                & \iff x\in A \cap (f^{-1}(B))^c\\
            \end{align*}
            
        \end{proof}
    \end{answer}

    \item Prove that $f$ is continuous on $A$ if and only if:
    \begin{equation}
        \label{eq:closed-preimage}
        \begin{gathered}
            \text{for every \emph{closed} subset }F\subseteq\mathbb{R}^m,\text{ then}\\
            f^{-1}(F)=A\cap E\text{ for some \emph{closed} subset }E\subseteq\mathbb{R}^n.
        \end{gathered}
    \end{equation}
    (Note that if the domain $A$ is all of $\mathbb{R}^n$, we can simplify
    this to $f^{-1}(F)$ is closed in $\mathbb{R}^n$ for every $F$ closed
    in $\mathbb{R}^m$.)
    \begin{answer}
        \begin{proof}
            Suppose \eqref{eq:closed-preimage} holds. We have to show $f$ is continuous on $A$.
            
            Let arbitrary $\open V\subseteq\mathbb{R}^m$. 
            Then $V^c$ is closed. 
            
            By \eqref{eq:closed-preimage}, $f^{-1}(V^c) = A\cap E$ for some $\closed E\subseteq\mathbb{R}^n$.
            
            By [a], $f^{-1}(V^c) = A \cap (f^{-1}(V))^c$
            \begin{align*}
                & A \cap E = A \cap (f^{-1}(V))^c \\
                & A \setminus E = A \setminus (f^{-1}(V))^c  & f^{-1}(V)\subseteq A\\
                & A \cap E^c = f^{-1}(V) \\
            \end{align*}

            Since $E^c$ and $V$ are open, by the open-set characterization of continuity, we have $f$ continuous on $A$.
            \vspace{1em}

            Now, suppose $f$ continuous on $A$. We have to show \eqref{eq:closed-preimage}.

            By the open-set characterization of continuity,
            we can take arbitrary $\open V\subseteq\mathbb{R}^m$ and let $\open U\subseteq\mathbb{R}^n$ such that 
            \[f^{-1}(V) = A\cap U\]

            By [a], $f^{-1}(V^c) = A \cap (f^{-1}(V))^c$
            \begin{align*}
                & f^{-1}(V^c) = A \cap (A\cap U)^c \\
                & f^{-1}(V^c) = A \cap (A^c\cup U^c) \\
                & f^{-1}(V^c) = (A \cap A^c)\cup (A\cap U^c) \\
                & f^{-1}(V^c) = (A\cap U^c) \\
            \end{align*}
            Since $U$ and $V$ are open, their complements are closed so we have shown \eqref{eq:closed-preimage} for arbitrary sets.
        \end{proof}
    \end{answer}

    \item Show that $B=\{(x,y)\in\mathbb{R}^2:xy=1\}$ is a closed subset
    of $\mathbb{R}^2$. \emph{Hint:} Consider $g(x,y)=xy$.
    \begin{answer}
        \begin{lemma}
            Let $g:\mathbb{R}^2\to\mathbb{R}$ such that $g(x,y) = xy$. $g$ is continuous on $\mathbb{R}^2$.

            Let arbitrary $p = (p_1,p_2)\in\mathbb{R}^2$. Take a sequence $(s_k)_{k\in\mathbb{N}} = (x_k, y_k)_{k\in\mathbb{N}}$ that converges to $p$.

            Therefore, $(x_k)\to p_1$ and $(y_k)\to p_2$
            \begin{align*}
                & g(x_k,y_k) = x_ky_k \\
                & \lim_{k\to\infty} g(x_k,y_k) = \lim_{k\to\infty}x_ky_k \\
                & \lim_{k\to\infty} g(x_k,y_k) = p_1p_2 \\
                & \lim_{k\to\infty} g(x_k,y_k) = g(p) \\
            \end{align*}
            So $(g(s_k))$ converges to $g(p)$. Therefore, $g$ is continuous at $p$. So $g$ is continuous on $\mathbb{R}^2$.
        \end{lemma}

        \begin{lemma}
            $\{1\}\subseteq\mathbb{R}$ is closed.
            
            Let arbitrary $x\in\mathbb{R}\setminus\{1\}$. Take $\varepsilon = \abs{x - 1} > 0$. 
            Then, $B_\varepsilon(x)\subseteq\mathbb{R}\setminus\{1\}$.So $\mathbb{R}\setminus\{1\}$ is open. So $\{1\}\subseteq\mathbb{R}$ is closed.
        \end{lemma}
        
        \begin{proof}
            Let $g:\mathbb{R}^2\to\mathbb{R}$ such that $g(x,y) = xy$.
            Since $g$ is continuous and $\{1\}$ is closed in $\mathbb{R}$, part [b] implies that 
            $B = g^{-1}(\{1\})$ is closed in $\mathbb{R}^2$.
        \end{proof}
    \end{answer}
\end{parts}
\end{problem}

\begin{problem}
Let $f:A\subseteq\mathbb{R}^n\to\mathbb{R}^m$ be a map. For
$\mathbf{x}\in A$, we have
\[
    f(\mathbf{x})=\bigl(f_1(\mathbf{x}),\ldots,f_m(\mathbf{x})\bigr),
\]
where $f_j:A\subseteq\mathbb{R}^n\to\mathbb{R}$ is the $j^{\text{th}}$
\emph{component function} of $f$, for $1\leq j\leq m$.
\begin{parts}
    \item Let $\mathbf{x}_0\in A$. Prove using the sequential characterization
    of continuity that $f$ is continuous at $\mathbf{x}_0$ if and only if
    each $f_j$ is continuous at $\mathbf{x}_0$ for all $1\leq j\leq m$.
    \begin{answer}
        \begin{proof}
            \begin{align*}
            &f \text{ is continuous at } x_0\\
            &\iff \forall (x_k)_{k\in\mathbb{N}}\subseteq A:\ \lim_{k\to\infty}x_k = x_0 \implies \lim_{k\to\infty}f(x_k) = f(x_0)\\
            &\iff \forall (x_k)_{k\in\mathbb{N}}\subseteq A:\ \lim_{k\to\infty}x_k = x_0 \implies\forall j\in[m]:\ \lim_{k\to\infty}f_j(x_k) = f_j(x_0)\\
            &\iff\forall j\in[m]:\ \forall (x_k)_{k\in\mathbb{N}}\subseteq A:\ \lim_{k\to\infty}x_k = x_0 \implies \lim_{k\to\infty}f_j(x_k) = f_j(x_0)\\
            &\iff\forall j\in[m]:\ f_j \text{ is continuous at }x_0\\
            \end{align*}
            We've shown the step going from line $2\to3$ of this in $A2$.
        \end{proof}
    \end{answer}

    \item Let $\pi_j:\mathbb{R}^m\to\mathbb{R}$ be the map
    $\pi_j(y_1,\ldots,y_m)=y_j$, the projection onto the $j^{\text{th}}$
    component. We proved in class directly using the definition of continuity
    that $\pi_j$ is continuous on $\mathbb{R}^m$. Let
    $\iota_j:\mathbb{R}\to\mathbb{R}^m$ be the map
    $\iota_j(t)=(0,\ldots,0,t,0,\ldots,0)$ where the $t$ is in the
    $j^{\text{th}}$ position. Prove directly using the definition of
    continuity that $\iota_j$ is continuous on $\mathbb{R}$.
    \begin{answer}
        \begin{proof}
            Let arbitrary $p\in\mathbb{R}$. Let arbitrary $\varepsilon > 0$. Take $\delta = \varepsilon$.

            We can see that
            \[\norm{\iota_j(x) - \iota_j(p)} = \norm{(0,\dots,0,x-p,0,\dots,0)} = \abs{x-p}\]
            Now, \[\forall x\in\mathbb{R}:\ \abs{x-p} < \delta = \varepsilon \implies \norm{\iota_j(x) - \iota_j(p)} = \abs{x-p} < \varepsilon\]
            So $\iota_j$ is continuous on $\mathbb{R}$
        \end{proof}
    \end{answer}

    \item Find a way to use the maps $\pi_j$, $\iota_j$ in [b] to give another
    proof that $f$ is continuous on $A$ if and only if each $f_j$ is
    continuous on $A$ for $1\leq j\leq m$.
    \begin{answer}
        \begin{proof}
            We can easily see that
            \begin{equation}
                \label{eq:projection-decomposition}
                \begin{gathered}
                    \forall j\in[m]:\ (\pi_j\circ f)(x) = f_j(x)
                \end{gathered}
            \end{equation}
            And that
            \begin{equation}
                \label{eq:component-composition}
                \begin{gathered}
                    f(x) = \sum_{j=1}^{m}(\iota_j\circ f_j)(x)
                \end{gathered}
            \end{equation}
            We also know from class that 
            \begin{itemize}
                \item composition of continuous maps are continuous;
                \item liniar combinations of continuous maps are continuous;
            \end{itemize}

            First, assume $f$ is continuous on $A$. We also have that $\pi_j$ is continuous on $\mathbb{R}^m$.
            Therefore, $(\pi_j\circ f)$ is continuous on $A$. So, by \eqref{eq:projection-decomposition}, $\forall j\in[m]:\ f_j$ continuous on $A$.
            
            Secondly, assume $f_j$ is continuous on $A$ for all $j\in[m]$. $\iota_j$ is also continuous on $\mathbb{R}$,
            so $(\iota_j\circ f_j)$ is continuous on $A$ for all $j\in[m]$.

            $\sum_{j=1}^{m}(\iota_j\circ f_j)(x)$ is a liniar combination of continuous maps so it's also continuous. 
            Therefore, by \eqref{eq:component-composition}, $f$ is continuous on $A$.
        \end{proof}
    \end{answer}
\end{parts}
\end{problem}

\begin{problem}
Let $\mathbb{R}^{n^2}$ denote the space of $n\times n$ (square) matrices.
Recall we can multiply two $n\times n$ matrices $A,B$ to obtain a matrix $AB$.
\begin{parts}
    \item Let $A,B\in\mathbb{R}^{n^2}$. Recall that the Euclidean inner
    product on this space is
    $\langle A,B\rangle=\sum_{i=1}^n\sum_{j=1}^n A_{ij}B_{ij}$.
    Prove that $\langle A,B\rangle=\operatorname{tr}(A^T B)$, where $A^T$
    is the transpose of $A$, and $\operatorname{tr}(C)$ is the trace of $C$.
    \begin{answer}
        \begin{proof}
            \begin{align*}
                & (A^TB)_{ii} = \sum_{j=1}^{n}A^T_{ij}B_{ji} = \sum_{j=1}^{n}A_{ji}B_{ji} \\
                & tr(A^TB) = \sum_{i=1}^{n}(A^TB)_{ii}  = \sum_{i=1}^{n} \sum_{j=1}^{n}A_{ji}B_{ji} = \ip{A}{B}
            \end{align*}
        \end{proof}
    \end{answer}

    \item Consider the map $g:\mathbb{R}^{n^2}\to\mathbb{R}^{n^2}$ given by
    $g(A)=A^T A$. Prove that $g$ is continuous on $\mathbb{R}^{n^2}$.
    \begin{answer}
        \begin{proof}
            Let arbitrary $M\in\mathbb{R}^{n^2}$. Let arbitrary sequence $(A_k)_{k\in\mathbb{N}}$ in $\mathbb{R}^{n^2}$ such that
            \[\lim_{k\to\infty}A_k = M\]
            Now, we've shown previously that for $\mathbb{R}^n$ convergence must happen component-wise, and it is trivial that this is also true in $\mathbb{R}^{n^2}$.
            So:
            \[\lim_{k\to\infty}A_k = M \iff \forall i,j\in[n]:\ \lim_{k\to\infty}(A_k)_{ij} = M_{ij}\]

            Now consider $(g(A_k))_{ij}$:
            \begin{align*}
                (g(A_k))_{ij} = (A_k^TA_k)_{ij} & = \sum_{r=1}^{n}(A_k)_{ri}(A_k)_{rj}\\
                \lim_{k\to\infty} (g(A_k))_{ij} & = \sum_{r=1}^{n}\lim_{k\to\infty}(A_k)_{ri}(A_k)_{rj} \\ 
                & = \sum_{r=1}^{n}M_{ri}M_{rj}\\
                & = (g(M))_{ij}
            \end{align*}
            And this holds $\forall i,j\in[n]$, so
            \[\lim_{k\to\infty}g(A_k) = g(M)\]
            Therefore, $g$ is continuous at $M$ and since $M$ was arbitrary, $g$ is continuous on $\mathbb{R}^{n^2}$.
        \end{proof}
    \end{answer}

    \item Prove that the subset of $n\times n$ \emph{orthogonal} matrices
    is a \emph{compact} subset of $\mathbb{R}^{n^2}$.
    \begin{answer}
        \begin{proof}
            Let \[\mathbb{E} = \{A\in\mathbb{R}^{n^2}:\ A^T = A^{-1}\}\]
            Equivalently,
            \[A\in\mathbb{E}\iff A^TA = I\]
            Therefore,
            \[\mathbb{E} = g^{-1}(\{I_{n}\})\]
            Since $g$ is continuous and $\{I_n\}$ is a closed singleton set, from $1[b]$ we get that $E$ is also closed.

            Now take arbitrary $A\in\mathbb{E}$.
            \begin{align*}
                & A^TA = I_n \\
                & tr(A^TA) = tr(I_n) = n\\
                & \ip{A}{A} = n\\
                & \norm{A} = \sqrt{n} 
            \end{align*}
            So $\mathbb{E}$ is bounded by $\sqrt{n}$.

            Therefore, $\mathbb{E}$ is closed and bounded, so it is compact.
        \end{proof}
    \end{answer}
\end{parts}
\end{problem}

\begin{problem}
Let $A\subseteq\mathbb{R}^n$ be nonempty, and fix a point
$\mathbf{x}\in\mathbb{R}^n$.
\begin{parts}
    \item If $A$ is \emph{compact}, show that there exists at least one point
    $\mathbf{y}\in A$ such that
    $\lVert\mathbf{y}-\mathbf{x}\rVert\geq\lVert\mathbf{a}-\mathbf{x}\rVert$
    for all $\mathbf{a}\in A$. That is, $\mathbf{y}$ is a point in $A$ that
    is \emph{furthest} from $\mathbf{x}$.

    \emph{Hint:} Find a way to use the Extreme Value Theorem.
    \begin{answer}
        Consider the map $f:A\subseteq\mathbb{R}^n\to\mathbb{R}$ given by 
        \[f(a) = \norm{a-x} = \sqrt{\sum_{i=1}^{n}(a_i-x_i)^2}\]
        We want to show that $f$ is continuous.

        Let arbitrary $a\in A$, let arbitrary $\varepsilon > 0$. Choose $\delta = \varepsilon$.
        
        Take arbitrary $b\in A$ such that $\norm{b-a} < \delta$.

        From the triangle inequality, we get that
        \[\abs{\norm{b-x} - \norm{a-x}} \leq \norm{b-a}\]
        So
        \[\abs{\norm{b-x} - \norm{a-x}} \leq \norm{b-a} < \delta = \varepsilon\]
        \[\abs{\norm{b-x} - \norm{a-x}} < \varepsilon\]
        \[\abs{f(b) - f(a)} < \varepsilon\]

        So $f$ is continuous on $A$ and $A$ is compact. By the Extreme Value Theorem,
        \[\exists y\in A:\ \forall a\in A:\  f(y) \geq f(a)\]
        So
        \[\exists y\in A:\ \forall a\in A:\  \norm{y-x} \geq \norm{a-x}\]
    \end{answer}

    \item Give an example to show that $\mathbf{y}$ is not unique in general.
    (There are many possible answers.)
    \begin{answer}
        Let $A = [-1,1]$ and $x = 0$. Then $y$ can be both $-1$ and $1$ as both achieve the max distance from $x$.

        Also we should note that $A$ is compact because it is a closed interval in $\mathbb{R}$. 
    \end{answer}

    \item If $A$ is \emph{not compact}, there may not exist any such
    $\mathbf{y}$. If $A$ is noncompact because it is unbounded, then this
    is trivial to see, as there are points in $A$ arbitrarily far from the
    origin, and hence arbitrarily far from any $\mathbf{x}$. More interesting
    is if $A$ is noncompact because it is bounded but \emph{not closed}.
    Give an example of such an $A$ and $\mathbf{x}$ such that there is no
    farthest point in $A$ from $\mathbf{x}$. (There are many possible answers.)
    \begin{answer}
        Let $A = (0,1)$ and $x = 0$. In this case, there is no $y$ that achieves the maximum distance away from $x$ as you can always get closer to $1$.
    \end{answer}
\end{parts}
\end{problem}

\begin{problem}
Let $f:\mathbb{R}^n\to\mathbb{R}$ be continuous on $\mathbb{R}^n$.
Suppose that $f$ has the following property: for any $\varepsilon>0$,
there exists $M>0$ such that if $\lVert\mathbf{x}\rVert>M$, then
$|f(\mathbf{x})|<\varepsilon$. (We say that $f$ \emph{goes to zero at infinity}.)
\begin{parts}
    \item Show that $f$ is bounded on $\mathbb{R}^n$.
    \begin{answer}
        \begin{proof}
            Take $\varepsilon = 1$. There exists $M > 0$ such that 
            \[\norm{x} > M \implies \abs{f(x)} < 1 \]
            Now consider the closed ball $\overline{B_M(0)}$.
            \[x\in \overline{B_M(0)} \iff \norm{x} \leq M \]
            $f$ is continuous and $\overline{B_M(0)}$ is closed and bounded, so $f(\overline{B_M(0)})$ is also bounded.
            
            Let $K$ be the bound of $f(\overline{B_M(0)})$.

            Now, if $\norm{x} \leq M$, then $f(x)$ bounded by $K$ and if $\norm{x} > M$, $f(x)$ bounded by $\varepsilon = 1$.
            
            Therefore, $f$ bounded by $\max K,1$.
        \end{proof}
    \end{answer}

    \item Suppose there exists some point $\mathbf{x}_0\in\mathbb{R}^n$
    such that $f(\mathbf{x}_0)>0$. Prove that there exists a point
    $\tilde{\mathbf{x}}$ such that $f(\mathbf{x})\leq f(\tilde{\mathbf{x}})$
    for all $\mathbf{x}\in\mathbb{R}^n$. (That is, show that $f$
    \emph{attains a global maximum} on $\mathbb{R}^n$.)

    [Note that we cannot \emph{directly} use the Extreme Value Theorem
    because $\mathbb{R}^n$ is not compact. Find a creative way to use it
    nevertheless.]
    \begin{answer}
        \begin{proof}
            Take $\varepsilon = f(x_0) > 0$. There exists $M > 0$ such that
            \[\norm{x} > M \implies \abs{f(x)} < f(x_0)\]

            Therefore, for any $x$ ouside the closed ball of size $M$ centered at $0$, $f(x)$ is bounded by $f(x_0)$.

            For all $x$ inside the closed ball, we can apply the EVT on that subset to show that there exists $\bar{x}$ for which $f(x) \leq f(\bar{x})$

            Now simply take $\tilde{x}\in\{x_0, \bar{x}\}$ for which $f(\tilde{x}) = \max\{f(x_0),f(\bar{x})\}$.
        \end{proof}
    \end{answer}

    \item Show that the assumption that there exists
    $\mathbf{x}_0\in\mathbb{R}^n$ with $f(\mathbf{x}_0)>0$ is necessary
    for part [b]. That is, find an example where $f$ \emph{goes to zero
    at infinity}, with $f(\mathbf{x})\leq0$ for all
    $\mathbf{x}\in\mathbb{R}^n$, and such that $f$ does not attain a global
    maximum on $\mathbb{R}^n$. [Take $n=1$ for simplicity.]
    \begin{answer}
        Let $f(x) = -e^{-x}$. $f(x)\leq 0$ for all $x$ and $\lim_{x\to\infty}f(x) = 0$.
        
        $f$ does not have a global maximum.
        \begin{proof}
            Assume for the sake of contradiction that $f$ has a global maximum. Then, 
            \[\exists \tilde{x}:\ \forall x\in\mathbb{R}:\ f(\tilde{x}) \geq f(x)\]
            However
            \begin{align*}
                &\tilde{x} + 1 > \tilde{x}\\
                &-(\tilde{x} + 1) < -(\tilde{x})\\
                &e^{-(\tilde{x} + 1)} < e^{-(\tilde{x})}\\
                &-e^{-(\tilde{x} + 1)} > -e^{-(\tilde{x})}\\
                &f(\tilde{x} + 1) > f(\tilde{x})\\
            \end{align*}

            Which is a contradiction.
        \end{proof}
    \end{answer}
\end{parts}
\end{problem}

\end{document}
